\begin{lemma}[An orthonormal basis adapted to a planar rotated frame]
Let $n\ge 2$, let $e,v\in\mathbb{R}^n$ satisfy $\lVert e\rVert=\lVert v\rVert=1$ and
$\langle e,v\rangle=0$, and let $\theta\in\mathbb{R}$. Put
\[
  e_r=\cos\theta\,e+\sin\theta\,v,
  \qquad
  e_\theta=-\sin\theta\,e+\cos\theta\,v .
\]
Then there are indices $i_0\ne i_1$ in $\{1,\dots,n\}$ and an orthonormal basis
$(b_i)_{i=1}^{n}$ of $\mathbb{R}^n$ with
\[
  b_{i_0}=e_r,\qquad b_{i_1}=e_\theta,
\]
and such that every remaining basis vector is orthogonal to the whole plane
$\Pi=\operatorname{span}\{e,v\}$: for every $i$ with $i\ne i_0$ and $i\ne i_1$,
\[
  \langle e,b_i\rangle=0
  \qquad\text{and}\qquad
  \langle v,b_i\rangle=0 .
\]
\end{lemma}
\begin{proof}
Throughout we use that $\langle e,e\rangle=\lVert e\rVert^2=1$,
$\langle v,v\rangle=\lVert v\rVert^2=1$, and $\langle e,v\rangle=\langle v,e\rangle=0$,
the last by symmetry of the real inner product, together with
$\cos^2\theta+\sin^2\theta=1$.

\emph{Step 1: $\{e_r,e_\theta\}$ is an orthonormal pair.} Expanding by
bilinearity,
\[
  \langle e_r,e_r\rangle
  =\cos^2\theta\,\langle e,e\rangle+2\cos\theta\sin\theta\,\langle e,v\rangle
    +\sin^2\theta\,\langle v,v\rangle
  =\cos^2\theta+\sin^2\theta=1 ,
\]
and likewise
\[
  \langle e_\theta,e_\theta\rangle
  =\sin^2\theta\,\langle e,e\rangle-2\sin\theta\cos\theta\,\langle e,v\rangle
    +\cos^2\theta\,\langle v,v\rangle
  =\sin^2\theta+\cos^2\theta=1 .
\]
Since a norm is nonnegative and $t\mapsto t^2$ is injective on $[0,\infty)$,
$\lVert e_r\rVert^2=1$ gives $\lVert e_r\rVert=1$, and similarly $\lVert e_\theta\rVert=1$.
Finally
\[
  \langle e_r,e_\theta\rangle
  =-\cos\theta\sin\theta\,\langle e,e\rangle+\cos^2\theta\,\langle e,v\rangle
    -\sin^2\theta\,\langle v,e\rangle+\sin\theta\cos\theta\,\langle v,v\rangle
  =-\cos\theta\sin\theta+\sin\theta\cos\theta=0 ,
\]
and $\langle e_\theta,e_r\rangle=0$ by symmetry.

\emph{Step 2: extension to an orthonormal basis.} Since $n\ge 2$ the index set
$\{1,\dots,n\}$ contains two distinct indices; fix $i_0\ne i_1$. Define a
family on the two-element subset $s=\{i_0,i_1\}$ by sending $i_0$ to $e_r$ and
$i_1$ to $e_\theta$. By Step 1 this family is orthonormal: each member is a unit
vector, and the two distinct members are mutually orthogonal. The index type has
exactly $n$ elements and $\dim\mathbb{R}^n=n$, so an orthonormal family indexed
by a subset of it extends to an orthonormal basis $(b_i)_{i=1}^{n}$ of
$\mathbb{R}^n$ agreeing with the given family on $s$. Hence $b_{i_0}=e_r$ and
$b_{i_1}=e_\theta$.

\emph{Step 3: the planar frame spans $\Pi$.} Inverting the rotation, we claim
\begin{equation}\label{eq:apf-dec}
  e=\cos\theta\,e_r+(-\sin\theta)\,e_\theta,
  \qquad
  v=\sin\theta\,e_r+\cos\theta\,e_\theta .
\end{equation}
Both are identities between linear combinations of $e$ and $v$, so it suffices
to compare coefficients. For the first, the coefficient of $e$ on the right is
$\cos\theta\cdot\cos\theta+(-\sin\theta)\cdot(-\sin\theta)
=\cos^2\theta+\sin^2\theta=1$, and the coefficient of $v$ is
$\cos\theta\sin\theta+(-\sin\theta)\cos\theta=0$. For the second, the
coefficient of $e$ is $\sin\theta\cos\theta+\cos\theta(-\sin\theta)=0$ and the
coefficient of $v$ is $\sin^2\theta+\cos^2\theta=1$. This proves
\eqref{eq:apf-dec}.

\emph{Step 4: the remaining basis vectors are orthogonal to $\Pi$.} Let $i$
satisfy $i\ne i_0$ and $i\ne i_1$. Since $(b_j)_j$ is orthonormal, distinct
indices give orthogonal vectors, so
\[
  \langle b_i,e_r\rangle=\langle b_i,b_{i_0}\rangle=0,
  \qquad
  \langle b_i,e_\theta\rangle=\langle b_i,b_{i_1}\rangle=0 ,
\]
using $b_{i_0}=e_r$ and $b_{i_1}=e_\theta$ from Step 2. Now pair $b_i$ with the
first identity of \eqref{eq:apf-dec} and use linearity of the inner product in
its second argument:
\[
  \langle b_i,e\rangle
  =\cos\theta\,\langle b_i,e_r\rangle+(-\sin\theta)\,\langle b_i,e_\theta\rangle
  =\cos\theta\cdot 0+(-\sin\theta)\cdot 0=0 ,
\]
and with the second identity,
\[
  \langle b_i,v\rangle
  =\sin\theta\,\langle b_i,e_r\rangle+\cos\theta\,\langle b_i,e_\theta\rangle
  =0 .
\]
By symmetry of the real inner product, $\langle e,b_i\rangle=\langle b_i,e\rangle=0$
and $\langle v,b_i\rangle=\langle b_i,v\rangle=0$. This is the last assertion,
and the lemma is proved.
\end{proof}
